%% ===========================================================================
%% 4.1 Prueba de concepto: arranque en frío y latencia          (Vicente)
%% Reemplaza el bloque entre "\subsection{Prueba de concepto..." y el
%% comentario de 4.2 en main.tex. Requiere en assets/images/poc-coldstart.pdf
%% la figura generada por poc/scripts/analizar.py (ya copiada).
%%
%% Presupuesto: ~350 palabras. Metodología ~110, interpretación ~140,
%% limitaciones ~80. La tabla y la figura no cuentan.
%%
%% Autoría (Anexo A): metodología, tabla y figura provienen del código y de
%% los datos (nivel 2, código auxiliar asistido). Los bloques marcados
%% [REDACTAR] son interpretación y limitaciones: se escriben a mano (6.1).
%% Los hechos disponibles van en comentarios; la prosa la pone el autor.
%% ===========================================================================

\subsection{Prueba de concepto: arranque en frío y latencia}
\label{sec:poc}

Se desplegó una función trivial idéntica, que devuelve el identificador de la
instancia que responde y su tiempo de vida, en tres plataformas de la
Tabla~\ref{tab:cuadro-comparativo}: AWS Lambda (microVM, us-east-1, 128~MB), Google
Cloud Run (contenedor, us-east4, 256~MiB) y Cloudflare Workers (V8 isolate, borde,
128~MB). Desde un cliente en Chile se corrieron 25 ciclos por plataforma en la misma
ventana (21-09-2026): una petición tras inducir el arranque en frío y cinco
calientes, cada una con conexión TCP/TLS nueva. El frío se induce con el evento del
ciclo de vida propio de cada modelo (cambio de configuración, publicación de versión,
terminación del proceso), de modo que la petición medida llega sin instancia caliente
disponible; la clasificación frío/caliente usa lo que devuelve la función, no la
latencia. Se reportan la latencia total percibida y la penalización de arranque (p50
frío menos p50 caliente), que cancela la red. Código, datos y notas en el Anexo~B.

\imagen{assets/images/poc-coldstart.pdf}{Latencia total desde el cliente en arranque en frío y en caliente, por plataforma (25 ciclos forzados, 150 peticiones por plataforma, 21-09-2026). La caja cubre el rango intercuartílico, la línea es la mediana y los puntos son valores atípicos.}{poc}

\tabla{Resumen de mediciones de la prueba de concepto (latencia total desde el cliente, ms)}{poc-resumen}{|p{3cm}|p{2cm}|p{2cm}|p{2.1cm}|p{2.1cm}|p{3.2cm}|}{%
    \filaTabla{\textbf{Plataforma} & \textbf{Frío p50} & \textbf{Frío p95} & \textbf{Caliente p50} & \textbf{Caliente p95} & \textbf{n frío/caliente, región, fecha}}
    \filaTabla{AWS Lambda & 838 & 890 & 442 & 507 & 25/125, us-east-1, 21-09-2026}
    \filaTabla{Cloudflare Workers & 569 & 591 & 567 & 574 & 107/43, borde (colo GIG), 21-09-2026}
    \filaTabla{Google Cloud Run & 922 & 1\,656 & 203 & 358 & 24/125, us-east4, 21-09-2026}
}

%% ---------------------------------------------------------------------------
%% [REDACTAR] Interpretación (~140 palabras, autoría propia).
%% Hechos disponibles, en el orden sugerido:
%%  1. Penalización p50: Lambda +395 ms, Workers +1 ms, Cloud Run +720 ms
%%     (results/resumen.csv). El orden sigue al tipo de aislamiento:
%%     isolate < microVM < contenedor. Es hallazgo propio.
%%  2. En Workers el arranque existe (107 isolates recién creados declararon
%%     first_request) pero el cliente no lo percibe: la latencia es la del
%%     viaje al borde. En Lambda los 25 fríos cayeron exactamente en la primera
%%     petición de cada ciclo.
%%  3. En caliente el orden se invierte y lo manda la red, no el modelo:
%%     Google termina la conexión TCP/TLS en un punto de presencia en Chile
%%     (TCP 16 ms) y reenvía por red interna a us-east4; Lambda conecta directo
%%     a Virginia (TCP ~140 ms). La latencia total incluye la arquitectura de
%%     front-end de cada proveedor.
%%  4. Cloud Run es la más dispersa (306-1682 ms): al terminar el proceso la
%%     plataforma empieza a reponer el contenedor por su cuenta y la petición
%%     paga la parte del arranque que falta cuando llega. El arranque del
%%     contenedor que registra la propia plataforma es parejo (~1,0 s).
%%  5. Cierre posible: la sección 2.2 desarrolla causas y mitigaciones; estas
%%     mediciones son la línea base sin mitigación (sin concurrencia
%%     aprovisionada, sin CPU boost, sin instancias mínimas).
%% ---------------------------------------------------------------------------
\textbf{[REDACTAR: interpretación]}

%% ---------------------------------------------------------------------------
%% [REDACTAR] Limitaciones (~80 palabras, autoría propia). Cuatro hechos:
%%  a. Cloud Run corre con 256 MiB porque Node.js 24 no cabe en 128 MiB en esa
%%     plataforma (el contenedor moría por memoria entre peticiones); Workers
%%     tiene 128 MB fijos y Lambda queda en 128 MB. En Lambda la CPU es
%%     proporcional a la memoria; en Cloud Run se fija aparte (1 vCPU).
%%  b. Los tres mecanismos para inducir el frío son distintos porque los tres
%%     ciclos de vida lo son; el de Cloud Run (terminar el proceso) es
%%     instrumentación de la prueba, no comportamiento de producción. Crear
%%     una revisión nueva no sirve: la plataforma la valida arrancando un
%%     contenedor antes de enrutar tráfico (24/25 primeras peticiones llegaron
%%     a un contenedor ya despierto en una corrida previa de la misma noche).
%%  c. La latencia al borde depende de la ruta del ISP: el 19-09 Workers
%%     caliente dio 231 ms y el 21-09, 567 ms, con el mismo colo (GIG). Dentro
%%     de una misma ventana la comparación entre plataformas sí es válida.
%%  d. Función trivial y un solo cliente: no representa cargas reales ni
%%     concurrencia.
%% ---------------------------------------------------------------------------
\textbf{[REDACTAR: limitaciones]}


%% ===========================================================================
%% ANEXO B: Código y mediciones de la prueba de concepto           (Vicente)
%% Reemplaza el bloque "\anexo{B}{...}" ... "\clearpage" de main.tex.
%% Los archivos de assets/code/ son copia de poc/ (rama PoC del repositorio).
%% Overleaf debe compilar con -shell-escape (minted).
%% ===========================================================================

\anexo{B}{Anexo B: Código y mediciones de la prueba de concepto}

\subsection*{B.1 Repositorio y trazabilidad}

El código, la configuración y los datos crudos están en el repositorio del grupo,
carpeta \texttt{poc/}, rama \texttt{PoC}. Corrida fuente de la sección~\ref{sec:poc}:
\texttt{data/mediciones\_20260921-0041\_forzado.csv} (450 filas, una por petición).
Corridas complementarias: \texttt{mediciones\_20260921-0005\_forzado.csv} (misma
noche, Cloud Run forzado por revisión nueva) y
\texttt{mediciones\_20260919-2148\_forzado.csv} (19-09, solo Lambda y Workers).
El resumen estadístico es \texttt{results/resumen.csv} y la latencia de arranque de
contenedor registrada por Google Cloud es \texttt{results/cloudrun\_startup.csv}.
El código auxiliar (funciones y scripts) se generó con asistencia de IA y fue
revisado y ejecutado por el autor; se declara en el Anexo~A.

\subsection*{B.2 Configuración de las plataformas}

\tabla{Configuración desplegada}{poc-config}{|p{2.6cm}|p{3.4cm}|p{3.4cm}|p{5cm}|}{%
    \filaTabla{\textbf{Plataforma} & \textbf{Región y aislamiento} & \textbf{Recursos} & \textbf{Evento que induce el frío}}
    \filaTabla{AWS Lambda & us-east-1; microVM & 128~MB (CPU proporcional), Node.js 24.x, Function URL & Cambio de la variable de entorno \texttt{POC\_MARKER}: el entorno de ejecución queda inválido y la invocación siguiente crea uno nuevo}
    \filaTabla{Cloudflare Workers & Borde (colo observado: GIG); V8 isolate & 128~MB fijos, plan Free & Publicación de una versión con \texttt{wrangler deploy}: las peticiones siguientes llegan a isolates nuevos (\texttt{first\_request})}
    \filaTabla{Google Cloud Run & us-east4; contenedor gen1 (gVisor) & 256~MiB, 1 vCPU, sin CPU boost, CPU solo durante solicitudes, min 0 / máx 1 instancia, concurrencia 1, Node.js 24 & \texttt{GET /salir}: el proceso responde y termina; con 0 instancias, la petición siguiente crea el contenedor}
}

\subsection*{B.3 Criterio de clasificación frío/caliente}

Cada petición se clasifica con lo que devuelve la función, nunca con su latencia:
en Lambda, frío si el \texttt{instance\_id} es distinto al de la petición anterior
(las peticiones secuenciales reutilizan el mismo entorno); en Workers, frío solo si
la instancia declara \texttt{first\_request}, porque Cloudflare reparte peticiones
consecutivas entre varios isolates vivos del mismo punto de presencia y el cambio
de identificador no implica arranque; en Cloud Run, frío si la edad del contenedor al
responder (\texttt{uptime\_ms}) es menor que la latencia medida de extremo a
extremo, es decir, si el contenedor nació dentro de la ventana de la petición. Este
último criterio no usa umbral: si el contenedor es más joven que la petición, solo
pudo nacer después de que ella salió del cliente.

\subsection*{B.4 Notas metodológicas sobre Cloud Run}

Con 128~MiB el contenedor excedía el límite de memoria (los registros del servicio
reportan entre 128 y 158~MiB usados) y era terminado entre peticiones, lo que
invalidaba la medición caliente; se subió a 256~MiB y la diferencia se declara.
Inducir el frío creando una revisión nueva no expone el arranque al cliente: la
plataforma valida la revisión arrancando un contenedor antes de enrutarle tráfico,
y en 24 de 25 ciclos la primera petición llegó a un contenedor ya despierto
(corrida de las 00:05). Por eso se induce con la terminación del proceso
(\texttt{GET /salir}, activo solo con la variable \texttt{POC\_SALIR\_HABILITADO=1}).
Tras la terminación se esperan 2~s: sin esa espera la petición llega a la instancia
que está muriendo (HTTP 503). La latencia de arranque de contenedor que registra la
propia plataforma (métrica \texttt{container/startup\_latencies}, ventana de la
corrida) fue de 26 arranques con media 1\,014~ms y p95 aproximado 1\,119~ms;
es otro instrumento (arranque visto por la plataforma, no latencia extremo a
extremo) y se reporta solo como contraste.

\subsection*{B.5 Código de las funciones}

\subsubsection*{AWS Lambda (\texttt{functions/lambda/index.mjs})}
\inputminted[fontsize=\scriptsize, breaklines]{javascript}{assets/code/lambda-index.mjs}

\subsubsection*{Cloudflare Workers (\texttt{functions/cloudflare/worker.js})}
\inputminted[fontsize=\scriptsize, breaklines]{javascript}{assets/code/cloudflare-worker.js}

\subsubsection*{Google Cloud Run (\texttt{functions/cloudrun/index.js})}
\inputminted[fontsize=\scriptsize, breaklines]{javascript}{assets/code/cloudrun-index.js}

\subsection*{B.6 Configuración y scripts de medición}

\subsubsection*{\texttt{config.json}}
\inputminted[fontsize=\scriptsize, breaklines]{json}{assets/code/config.json}

%% medir.py y analizar.py tienen varios cientos de líneas. Opción A: incluirlos
%% completos (descomentar; agrega ~12 páginas al anexo). Opción B: dejar solo la
%% referencia al repositorio (B.1) y este párrafo. Se recomienda B.
% \subsubsection*{\texttt{scripts/medir.py}}
% \inputminted[fontsize=\scriptsize, breaklines]{python}{assets/code/medir.py}
% \subsubsection*{\texttt{scripts/analizar.py}}
% \inputminted[fontsize=\scriptsize, breaklines]{python}{assets/code/analizar.py}
% \subsubsection*{\texttt{scripts/startup\_cloudrun.py}}
% \inputminted[fontsize=\scriptsize, breaklines]{python}{assets/code/startup_cloudrun.py}

El script de medición \texttt{scripts/medir.py} invoca cada URL por HTTP, mide la
latencia total, clasifica con el criterio de B.3 y escribe una fila por petición
con volcado inmediato; \texttt{scripts/analizar.py} produce
\texttt{results/resumen.csv}, las filas de la Tabla~\ref{tab:poc-resumen} y la
Figura~\ref{fig:poc}; \texttt{scripts/startup\_cloudrun.py} consulta la métrica de
arranque de contenedor de Cloud Monitoring para la ventana de la corrida. Ambos
scripts, con sus comentarios, están en el repositorio (B.1).

\clearpage
